\section{神经网络的函数空间视角}
\label{sec:14-Deep-Learning/10-Interpretability-Mathematics/01}

同一份代码、同一份数据、只换随机种子，训练出两个网络 A 和 B，测试准确率分别是 \(94.1\%\) 和 \(94.3\%\)，逐样本预测有九成以上一致。你想把它们融合一下，于是把权重逐元素平均，得到 \((\theta_A+\theta_B)/2\)。新模型的准确率是 \(10.2\%\)——十类分类的随机水平。

同一时间，另一位同事在做可解释性汇报：他把第三层某个通道的权重向量画出来，说这个神经元在检测``斜向边缘''。你追问换个种子重训之后哪个通道还是它，他答不上来。

这两件事共用一个根源。网络定义的是一个从输入到输出的函数，参数只是描述这个函数的一组坐标。坐标可以换，函数不变；坐标平均之后得到的，一般不是函数的平均。参数表里的数字对函数的行为有决定作用，但它们对函数的描述是高度冗余、而且不唯一的。深度模型把选择表示变成可优化的对象之后，得到的是一个函数——不是一份可以逐条朗读的参数说明书。

\subsection{参数到函数的映射不是单射}

把网络看成映射 \(\Phi:\theta\mapsto f_\theta\)，定义域是 \(\R^p\)，值域是某个函数集合。这个映射远不是一一对应，最容易验证的两类冗余如下。

一类是置换。把某个隐藏层的第 \(i\) 个与第 \(j\) 个单元互换位置，同时把下一层对应的输入权重也互换，网络计算出的函数逐点不变。宽度为 \(m_1,\dots,m_{L}\) 的多层网络因此至少有 \(\prod_l m_l!\) 个参数点给出同一个函数。宽度 512 的一层就贡献了 \(512!\) 这个量级的等价参数，远超宇宙中原子数目的估计。

另一类是连续的正缩放。ReLU 正齐次：\(\relu(ax)=a\,\relu(x)\) 对一切 \(a>0\) 成立。于是把某个单元的输入权重乘以 \(a\)、把它送往下一层的输出权重除以 \(a\)，函数同样不变：
\[
W_2\,\relu\bigl((aW_1)x\bigr)\cdot\frac1a=W_2\,\relu(W_1x).
\]
每个 ReLU 单元贡献一个连续参数 \(a\in(0,\infty)\)，这些自由度合起来构成一个正对角群，参数空间中给出一整条连通的等价轨道。条件对账：需要激活函数正齐次（ReLU 与 leaky ReLU 满足，sigmoid、tanh、GELU 都不满足）；带偏置时偏置须同步乘 \(a\)；训练目标中不能有权重范数惩罚，否则 \(L_2\) 项会在轨道上挑出范数最小的那个代表，函数不变而目标值变了。带归一化层的网络里也有类似的缩放冗余，上一单元讨论过它导致的``权重范数充当学习率''的效应。

于是参数空间上的等价关系是清楚的：\(\theta\sim\theta'\) 当且仅当 \(f_\theta=f_{\theta'}\)。被优化的对象实际生活在商空间 \(\R^p/\!\sim\) 上，而不是 \(\R^p\) 上。等价关系与商集的一般语言见\hyperref[sec:01-Mathematical-Language/02-Sets-Relations-Functions/02]{``关系与等价类''}。开头那次失败的权重平均现在有了解释：A 和 B 落在同一个函数附近的两条不同轨道上，把两组坐标做算术平均，结果落在轨道之间的某处，对应的函数与两者都无关。

\subsection{参数距离与函数距离互不控制}

映射 \(\Phi\) 不是单射，只是问题的一半。另一半是它两个方向上都不保距。

\begin{center}
\begin{tikzpicture}[>=stealth]
\draw[gray] (0,0) ellipse (2.5 and 1.9);
\node[below,font=\small] at (0,-2.1) {参数空间 \(\R^p\)};
\draw[gray] (7.6,0) ellipse (2.1 and 1.7);
\node[below,font=\small] at (7.6,-2.1) {函数空间};
\draw[dashed] plot[smooth] coordinates {(-1.5,0.55) (-0.7,1.05) (0.2,1.15) (1.1,0.75) (1.6,0.35)};
\foreach \p in {(-1.5,0.55),(-0.7,1.05),(0.2,1.15),(1.1,0.75),(1.6,0.35)} {
  \fill \p circle (1.8pt);
}
\node[below,font=\scriptsize] at (0.1,0.45) {同一函数的等价轨道};
\fill (6.9,0.85) circle (2.2pt);
\node[right,font=\scriptsize] at (7.05,0.95) {\(f\)};
\draw[->] (1.75,0.45) .. controls (3.6,1.5) .. (6.75,0.92);
\draw[->] (-1.35,0.75) .. controls (2.6,2.6) .. (6.78,1.02);
\fill (-1.1,-1.0) circle (1.8pt);
\fill (-0.75,-1.12) circle (1.8pt);
\node[above,font=\scriptsize] at (-0.9,-0.92) {参数改动看似很小};
\fill (6.6,-0.35) circle (2.2pt);
\fill (8.6,-1.05) circle (2.2pt);
\node[right,font=\scriptsize] at (8.7,-1.05) {\(g'\)};
\node[left,font=\scriptsize] at (6.5,-0.35) {\(g\)};
\draw[->] (-1.06,-0.99) .. controls (3.0,-0.9) .. (6.45,-0.42);
\draw[->] (-0.71,-1.14) .. controls (3.4,-2.3) .. (8.45,-1.05);
\end{tikzpicture}
\end{center}

\emph{一整条等价轨道可以压成函数空间里的一个点；反过来，若局部参数--函数 Jacobian 很大，数值上看似很小的参数改动也可能造成显著的行为差异。下半图表示放大效应，不表示映射不连续。}

图的上半部分是刚才那件事：一整条轨道被 \(\Phi\) 压成一个点，参数距离可以很大而函数距离为零。下半部分画的是放大效应：靠近输出层的一个权重只改千分之一，若它乘到的上游激活很大，所有样本的 logit 仍可能明显平移。这里必须保留一条边界：对固定的有限输入、有限权重和连续激活，\(\theta\mapsto f_\theta\) 本身是连续的；在有界区域上还能写出局部 Lipschitz 上界，所以``参数任意接近而函数仍隔着固定距离''并不成立。真正缺少的是一个不依赖数据尺度、网络深度与当前坐标的统一常数：仅报告参数欧氏距离，既不能证明行为接近，也不能证明行为不同。

要谈两个模型有多接近，度量应当定义在函数上。最常用的一个是相对数据分布的 \(L^2\) 距离
\[
\norm{f-g}_{L^2(P)}^2=\E_{x\sim P}\bigl[(f(x)-g(x))^2\bigr],
\]
在有限样本上用经验平均估计；分类任务里更合适的是预测一致率或输出分布的 KL 散度。这些量都依赖分布 \(P\)：两个函数在训练分布上几乎重合，在别的分布上可以完全不同，所以报告函数距离时必须一并报告在哪个分布上测的。函数作为空间中的点、以及由积分诱导的距离，见\hyperref[sec:12-Real-Analysis-Support/07-Lp-Spaces-and-Function-Geometry/01]{``\(L^p\) 范数把积分变成距离''}。

\subsection{坐标依赖的陈述解释不了模型}

把这条判断推到可解释性上，可以给出一个筛子：一个关于模型的陈述，如果在参数的等价变换下会改变，那它描述的是坐标而不是模型。

按这个筛子，``第三层第 17 个通道检测斜向边缘''不合格——置换一下单元编号，同一个函数就要换一种说法。``这个神经元的权重范数很大所以重要''也不合格，因为 ReLU 网络里可以在不改变函数的前提下把它的输入权重放大一百倍。``两个模型的权重余弦相似度是 0.02 所以学到了不同的东西''同样不合格，它没有排除置换与缩放。相比之下，``把这个通道的激活置零，测试集上第 3 类的召回率从 0.91 掉到 0.62''是合格的：它陈述的是对函数做一次干预之后行为的变化，置换单元编号不改变这个数值。

同一个筛子解释了为什么``看参数''这条路走不通，却没有否定所有内部分析。合格的做法都有一个共同形状：在函数或激活的层面上定义一个不随坐标改变的量，再把它测出来。方向而非单个坐标是常用的载体——若干单元的一个线性组合在等价变换下随之协变，用它定义的量可以是良定义的；对表示做无监督分解，找出一组与坐标无关的方向，也属于这条路线。这些方法把注意力从``哪个数字大''移到``函数在哪些方向上被区分''，换来这份良定义的是：得另外说明该方向是如何确定的，以及换一次训练还能不能找回它。

\subsection{函数视角改变的是提问方式}

回到最初的两个场景。权重平均失败不是 bug，它是在提醒你所处的空间不是线性的——若真要合并两个模型，需要先把它们对齐到同一条轨道上（近年的置换对齐工作正是在做这件事），或者干脆在函数层面合并，例如对输出取平均。至于神经元的语义，正确的问法不是``这个单元代表什么''，而是``对函数做哪种干预会让哪个行为发生可测的变化''。

这一节没有回答网络学到了什么，只是划出了提问的合法范围：陈述必须在参数的等价变换下不变。剩下的问题变成，训练过程究竟在函数空间里走了怎样一条路径。有一个极限情形下这条路径可以算得出来——网络宽度趋于无穷时，函数的演化退化成一个固定核上的线性过程。下一节从这个极限出发，看它解释了什么，又把深度模型最要紧的那一部分漏在了外面。
